% ═══════════════════════════════════════════════════════════════
% LEHRERHINWEISE: Elektromotor-Experiment (Kl. 9)
% ═══════════════════════════════════════════════════════════════

\documentclass[11pt, a4paper]{article}

\usepackage[utf8]{inputenc}
\usepackage[T1]{fontenc}
\usepackage[ngerman]{babel}

\usepackage{helvet}
\renewcommand{\familydefault}{\sfdefault}

\usepackage{microtype}
\usepackage[margin=2cm, top=2cm, bottom=2cm]{geometry}
\usepackage{enumitem}
\usepackage{tabularx}
\usepackage{booktabs}
\usepackage{array}
\usepackage{amssymb}

\usepackage{fancyhdr}
\usepackage{lastpage}

\usepackage{xcolor}
\usepackage{tcolorbox}
\tcbuselibrary{skins}

% ═══════════════════════════════════════════════════════════════
% FARBEN
% ═══════════════════════════════════════════════════════════════

\definecolor{physik}{HTML}{2c5aa0}
\definecolor{mittelgrau}{HTML}{888888}
\definecolor{erfolg}{HTML}{2E7D32}
\definecolor{warnung}{HTML}{E65100}

\colorlet{fachfarbe}{physik}

% ═══════════════════════════════════════════════════════════════
% KOPF- UND FUSSZEILE
% ═══════════════════════════════════════════════════════════════

\pagestyle{fancy}
\fancyhf{}
\renewcommand{\headrulewidth}{0.4pt}
\renewcommand{\footrulewidth}{0pt}
\lhead{\small Physik Klasse 9 | Elektromotor-Experiment}
\rhead{\small \textbf{LEHRERHINWEISE}}
\cfoot{\small\textcolor{mittelgrau}{Seite \thepage{} von \pageref{LastPage}}}

% ═══════════════════════════════════════════════════════════════
% BOXEN
% ═══════════════════════════════════════════════════════════════

\newtcolorbox{infobox}[1][]{
    colback=white,
    colframe=fachfarbe,
    coltitle=black,
    colbacktitle=white,
    boxrule=1pt,
    arc=2pt,
    left=6pt, right=6pt, top=4pt, bottom=4pt,
    before skip=8pt, after skip=8pt,
    fonttitle=\bfseries,
    title=#1
}

\newtcolorbox{warnbox}[1][]{
    colback=white,
    colframe=warnung,
    coltitle=black,
    colbacktitle=white,
    boxrule=1pt,
    arc=2pt,
    left=6pt, right=6pt, top=4pt, bottom=4pt,
    before skip=8pt, after skip=8pt,
    fonttitle=\bfseries,
    title=#1
}

% ═══════════════════════════════════════════════════════════════
% BEFEHLE
% ═══════════════════════════════════════════════════════════════

\setlength{\parindent}{0pt}
\setlength{\parskip}{6pt}

% ═══════════════════════════════════════════════════════════════
% DOKUMENT
% ═══════════════════════════════════════════════════════════════

\begin{document}

\begin{center}
    {\Large\bfseries\textcolor{fachfarbe}{LEHRERHINWEISE}}\\[0.3em]
    {\large Elektromotor-Experiment (40 Min)}
\end{center}
\hrule
\vspace{0.5em}

% ═══════════════════════════════════════════════════════════════
% ÜBERBLICK
% ═══════════════════════════════════════════════════════════════

\section*{Überblick}

\begin{tabularx}{\textwidth}{@{}lX@{}}
\textbf{Zielgruppe:} & Klasse 9 (nach Lernpfad LP-09a Gleichstrommotor) \\
\textbf{Dauer:} & 40 Minuten (35 Min Experiment + 5 Min Abbau) \\
\textbf{Sozialform:} & Partnerarbeit oder Kleingruppen (2--3 SuS) \\
\textbf{Material pro Tisch:} & 1 Motor, 1 Netzgerät, Kabel, 1 Stabmagnet, 1 Elektromagnet \\
\textbf{Arbeitsblätter:} & Experimentierkarte (1× pro Tisch), Forscherprotokoll (1× pro SuS) \\
\end{tabularx}

% ═══════════════════════════════════════════════════════════════
% ZEITPLAN
% ═══════════════════════════════════════════════════════════════

\section*{Zeitplan}

\begin{tabularx}{\textwidth}{@{}c|l|X@{}}
\toprule
\textbf{Min} & \textbf{Phase} & \textbf{Inhalt} \\
\midrule
0--3 & Einstieg & Material verteilen, Sicherheitshinweise, Aufbau-Checkliste \\
3--8 & Aufgabe 1--3 & Motor starten, Drehrichtung, Magnet-Vergleich \\
8--15 & Aufgabe 4--5 & Spannung verändern, Magnet-Abstand \\
15--22 & Aufgabe 6--7 & Motor stoppen, Lautstärke \\
22--32 & Knobelaufgaben & Totpunkt-Problem, Magnet-Bewertung (AFB III) \\
32--35 & Bonus/Fazit & Startspannung messen, Fazit notieren \\
35--40 & Abbau & Abbau-Checkliste, Material einsammeln \\
\bottomrule
\end{tabularx}

% ═══════════════════════════════════════════════════════════════
% VORBEREITUNG
% ═══════════════════════════════════════════════════════════════

\section*{Vorbereitung}

\begin{itemize}[noitemsep]
    \item Motoren vorab testen (funktionieren alle?)
    \item Netzgeräte auf 0\,V stellen
    \item Experimentierkarten laminieren (wiederverwendbar)
    \item Forscherprotokolle kopieren (1× pro SuS)
    \item Ggf. Nachschreiber separieren (schreiben MT-09a)
\end{itemize}

\begin{warnbox}[Typische Probleme]
\begin{itemize}[noitemsep]
    \item \textbf{Motor dreht nicht:} Kontakte prüfen, Magnet näher, Spannung erhöhen
    \item \textbf{Motor brummt nur:} Leicht anschieben (Totpunkt!), Spannung erhöhen
    \item \textbf{Elektromagnet zu schwach:} Mehr Windungen oder näher positionieren
    \item \textbf{Kabel heiß:} Kurzschluss! Sofort Spannung auf 0\,V
\end{itemize}
\end{warnbox}

% ═══════════════════════════════════════════════════════════════
% ERWARTUNGSHORIZONT
% ═══════════════════════════════════════════════════════════════

\section*{Erwartungshorizont}

\subsection*{Aufgabe 1: Motor starten (AFB I)}

\textbf{Erwartete Antwort:}
\begin{itemize}[noitemsep]
    \item Motor + Stromquelle (Netzgerät) + Kabel + Magnet (Stab- oder Elektromagnet)
    \item Akzeptabel: ``Strom und Magnetfeld''
\end{itemize}

\subsection*{Aufgabe 2: Drehrichtung ändern (AFB II)}

\textbf{Erwartete Antworten (2 von 3):}
\begin{enumerate}[noitemsep]
    \item Stromrichtung umpolen (Kabel am Netzgerät vertauschen)
    \item Magnetfeld umpolen (Stabmagnet umdrehen / Elektromagnet umpolen)
    \item \textit{Nicht:} Beides gleichzeitig (hebt sich auf!)
\end{enumerate}

\subsection*{Aufgabe 3: Magnet-Vergleich (AFB II)}

\textbf{Erwartete Antwort:}
\begin{itemize}[noitemsep]
    \item Beide funktionieren: ja / ja
    \item Unterschied: Elektromagnet oft schwächer, Feldstärke regelbar, braucht eigene Stromversorgung
\end{itemize}

\subsection*{Aufgabe 4: Spannung verändern (AFB I)}

\textbf{Erwartete Antwort:}
\begin{itemize}[noitemsep]
    \item Höhere Spannung $\rightarrow$ schnellere Drehung / höhere Drehzahl
    \item Akzeptabel: ``dreht schneller'', ``wird lauter''
\end{itemize}

\subsection*{Aufgabe 5: Magnet-Abstand (AFB I)}

\textbf{Erwartete Antwort:}
\begin{itemize}[noitemsep]
    \item Größerer Abstand $\rightarrow$ schwächeres Magnetfeld $\rightarrow$ langsamere Drehung
    \item Bei zu großem Abstand: Motor stoppt
\end{itemize}

\subsection*{Aufgabe 6: Motor stoppen (AFB II)}

\textbf{Erwartete Antworten (mind. 1):}
\begin{enumerate}[noitemsep]
    \item Magnet entfernen (kein Magnetfeld)
    \item Magnet weit weg (Feld zu schwach)
    \item Rotor im Totpunkt festhalten
    \item Spannung sehr niedrig (Reibung überwiegt)
\end{enumerate}

\subsection*{Aufgabe 7: Lautstärke (AFB II)}

\textbf{Erwartete Antwort:}
\begin{itemize}[noitemsep]
    \item Spannung beeinflusst Lautstärke (mehr Spannung = lauter)
    \item Untergrund beeinflusst Lautstärke (Resonanz)
    \item Magnet-Position kann Vibrationen beeinflussen
\end{itemize}

\subsection*{Knobelaufgabe A: Totpunkt-Problem (AFB III)}

\textbf{Erwartete Erklärung:}
\begin{itemize}[noitemsep]
    \item Im Totpunkt steht die Leiterschleife parallel zum Magnetfeld
    \item Die Lorentzkraft wirkt dann nicht tangential, sondern radial (nach innen/außen)
    \item $\Rightarrow$ Kein Drehmoment $\Rightarrow$ Motor bleibt stehen
\end{itemize}

\textbf{Skizze:} Rotor in horizontaler Position (0° oder 180°), Kraftpfeile zeigen nach innen/außen statt tangential.

\subsection*{Knobelaufgabe B: Magnet-Bewertung (AFB III)}

\textbf{Erwartete Antwort:}

\begin{tabularx}{\textwidth}{@{}l|X|X@{}}
& \textbf{Stabmagnet} & \textbf{Elektromagnet} \\
\midrule
\textbf{Vorteil} & Kein Strom nötig, kompakt, günstig, leicht & Feldstärke regelbar, abschaltbar \\
\textbf{Nachteil} & Feldstärke fix, kann entmagnetisieren & Braucht Stromversorgung, schwerer, teurer \\
\end{tabularx}

\textbf{Empfehlung:} Stabmagnet (für Spielzeugmotor: einfacher, leichter, billiger, zuverlässiger)

\subsection*{Bonus: Startspannung}

\textbf{Typische Werte:} 1--3\,V (abhängig von Motor und Reibung)

% ═══════════════════════════════════════════════════════════════
% DIFFERENZIERUNG
% ═══════════════════════════════════════════════════════════════

\section*{Differenzierung}

\begin{tabularx}{\textwidth}{@{}l|X@{}}
\toprule
\textbf{Niveau} & \textbf{Erwartung} \\
\midrule
G-Niveau (★) & Aufgaben 1, 4, 5 vollständig; Aufgaben 2, 3, 6, 7 teilweise \\
E-Niveau (★★) & Alle Aufgaben 1--7 vollständig \\
E+ (★★★) & Zusätzlich Knobelaufgaben A und B mit Begründung \\
\bottomrule
\end{tabularx}

% ═══════════════════════════════════════════════════════════════
% STOLPERSTEINE
% ═══════════════════════════════════════════════════════════════

\section*{Typische Schülerfehler}

\begin{itemize}[noitemsep]
    \item \textbf{Aufgabe 2:} ``Beides gleichzeitig umpolen'' als Lösung $\rightarrow$ Drehrichtung bleibt gleich!
    \item \textbf{Aufgabe 6:} ``Strom abschalten'' $\rightarrow$ Aufgabe verlangt ``obwohl Strom fließt''
    \item \textbf{Knobel A:} Verwechslung Totpunkt mit ``Motor ist kaputt''
    \item \textbf{Knobel B:} Nur Vor- oder Nachteile, keine Abwägung
\end{itemize}

% ═══════════════════════════════════════════════════════════════
% SICHERHEIT
% ═══════════════════════════════════════════════════════════════

\begin{warnbox}[Sicherheitshinweise]
\begin{itemize}[noitemsep]
    \item Maximale Spannung: 12\,V (Kleinspannung, ungefährlich)
    \item Kabel nicht unter Spannung umstecken
    \item Motor kann bei längerem Betrieb heiß werden
    \item Bei Kurzschluss: Sofort Netzgerät auf 0\,V
\end{itemize}
\end{warnbox}

\vfill

\begin{center}
\small\textcolor{mittelgrau}{Lehrerhinweise zum Forscherprotokoll | Physik Klasse 9}
\end{center}

\end{document}
